\documentclass[11pt]{article}
\usepackage[margin=1in]{geometry}
\usepackage{amsmath,amssymb}
\usepackage{bm}
\usepackage{booktabs}
\usepackage{hyperref}
\usepackage{xcolor}
\usepackage{longtable}
\usepackage{listings}
\usepackage{tcolorbox}
\usepackage{enumitem}

\hypersetup{colorlinks=true, linkcolor=blue!60!black, urlcolor=blue!60!black}

\lstset{basicstyle=\ttfamily\small, breaklines=true, frame=single,
  keywordstyle=\color{blue!70!black}, commentstyle=\color{green!50!black},
  backgroundcolor=\color{gray!5}}

\newcommand{\code}[1]{\texttt{#1}}
\newcommand{\placeholder}[1]{\textcolor{red!70!black}{\texttt{[#1]}}}

\setlist[itemize]{nosep, left=0pt}
\setlist[enumerate]{nosep, left=0pt}

\title{\textbf{PRISM v3}\\[4pt]
\large Pipeline Architecture and Model Documentation}
\author{Benjamin Coriat\\McGill University --- ECSE 551}
\date{April 2026}

\begin{document}
\maketitle

\begin{abstract}
Technical reference for the PRISM v3 codebase.
Every section maps to a code cell; every hyperparameter is listed with its default.
This document is the ``how''; for the ``why,'' see \textit{PRISM: Mathematical Foundations}.
\end{abstract}

\newpage
\tableofcontents
\newpage

%======================================================================
\section{Pipeline at a Glance}
%======================================================================

PRISM v3 lives in a single file (\code{prism\_v3.py}), fifteen cells
running end-to-end from raw Yahoo Finance data to a fully diagnosed backtest.
The flow is linear:

\begin{center}
\renewcommand{\arraystretch}{1.15}
\begin{tabular}{@{}r@{\;\;}l@{\qquad}r@{\;\;}l@{}}
\toprule
\textbf{Cell} & \textbf{What it does} & \textbf{Cell} & \textbf{What it does}\\
\midrule
1  & Install deps          & 9  & Training loop\\
2  & Imports \& logging    & 10 & Predict + error tracking\\
3  & Config dataclass      & 11 & Allocation + OT\\
4  & Order book tracker    & 12 & Backtest engine\\
5  & Download \& clean     & 13 & Vol-fit + metrics\\
6  & Weekly resample + GK  & 14 & 31 diagnostic plots\\
7  & Feature engineering   & 15 & Summary dump\\
8  & Stacked ensemble      &    & \\
\bottomrule
\end{tabular}
\end{center}

Everything downstream of cell~3 is parameterised by a single
\code{@dataclass Config} (Section~\ref{sec:config}).
No magic numbers live outside it.

%======================================================================
\section{Configuration}\label{sec:config}
%======================================================================

\begin{longtable}{@{}llp{6.8cm}@{}}
\toprule
\textbf{Parameter} & \textbf{Default} & \textbf{Role}\\
\midrule
\endhead
\multicolumn{3}{@{}l}{\textit{Date boundaries}}\\[2pt]
\code{START\_DATE}       & 2003-01-01 & Download start\\
\code{END\_DATE}         & 2025-12-31 & Download end\\
\code{TRAIN\_END}        & 2011-12-31 & Last in-sample date\\
\code{TEST\_START}       & 2012-01-01 & First OOS date\\[6pt]
\multicolumn{3}{@{}l}{\textit{Feature construction}}\\[2pt]
\code{WEEKLY\_FREQ}      & W-FRI     & Resample anchor\\
\code{CORR\_WINDOW}      & 26        & Rolling correlation window (weeks)\\
\code{RETURN\_LAGS}      & [1,2,4,8,12] & AR lag set\\
\code{VOL\_WINDOW}       & 8         & Return-vol rolling window\\
\code{VOL\_HISTORY}      & 26        & Depth $k$ of the $V$ matrix\\[6pt]
\multicolumn{3}{@{}l}{\textit{Model training}}\\[2pt]
\code{EWMA\_DECAY}       & 0.94      & Error-history decay\\
\code{TUNING\_TRIALS}    & 10        & Optuna trials per model\\
\code{TUNING\_TIMEOUT}   & 20\,s     & Optuna wall-clock cap\\
\code{CV\_SPLITS}        & 5         & \code{TimeSeriesSplit} folds\\
\code{TORCH\_EPOCHS}     & 100       & MLP max epochs\\
\code{TORCH\_PATIENCE}   & 12        & MLP early-stop patience\\
\code{TORCH\_BATCH}      & 64        & MLP batch size\\[6pt]
\multicolumn{3}{@{}l}{\textit{Allocation}}\\[2pt]
\code{LONG\_ONLY}        & False     & If true, $w_i\ge 0$\\
\code{W\_MAX}            & 0.20      & Single-asset cap\\
\code{L\_MAX}            & 1.6       & Gross leverage cap\\
\code{SIGMA\_TARGET}     & 0.15      & Ann.\ vol target\\
\code{RF\_ANNUAL}        & 0.045     & Risk-free rate\\[6pt]
\multicolumn{3}{@{}l}{\textit{Execution}}\\[2pt]
\code{TC\_VOL\_SENSITIVITY} & 0.05   & Half-spread vol sensitivity ($\alpha_{\mathrm{tc}}$)\\
\code{TC\_IMPACT\_SCALE}    & 0.02   & Market impact scale ($\beta_{\mathrm{tc}}$)\\
\code{SINKHORN\_EPS}     & 0.05      & OT regularisation scale\\
\code{SINKHORN\_ITER}    & 200       & Max Sinkhorn iterations\\
\code{SINKHORN\_TOL}     & $10^{-8}$ & Convergence threshold\\
\code{VOL\_LOOKBACK}     & 26        & Vol-target lookback\\
\code{VOL\_CAP}          & 1.0       & Max vol-target leverage (no leverage)\\
\bottomrule
\end{longtable}

%======================================================================
\section{Asset Universe}\label{sec:tickers}
%======================================================================

24 tickers covering global equity indices and commodities.
Any ticker with fewer than 100 daily observations after cleaning is dropped,
so $n$ is determined at runtime.

\begin{center}
\renewcommand{\arraystretch}{1.1}
\begin{tabular}{@{}llp{9.5cm}@{}}
\toprule
\textbf{Class} & \textbf{Region} & \textbf{Tickers}\\
\midrule
Equity & N.\ America  & \code{\^{}GSPC}, \code{\^{}NDX}, \code{\^{}DJI}, \code{\^{}RUT}, \code{\^{}GSPTSE}\\
       & Europe       & \code{\^{}FTSE}, \code{\^{}GDAXI}, \code{\^{}STOXX50E}, \code{\^{}FCHI}, \code{\^{}IBEX}\\
       & Asia-Pac     & \code{\^{}N225}, \code{\^{}HSI}, \code{\^{}AXJO}, \code{\^{}BSESN}, \code{\^{}NSEI}, \code{\^{}KS11}, \code{\^{}TWII}\\
       & LatAm        & \code{\^{}BVSP}, \code{\^{}MXX}\\
Vol    & N.\ America  & \code{\^{}VIX}\\
Commodity & Global   & \code{GC=F}, \code{SI=F}, \code{CL=F}, \code{NG=F}\\
\bottomrule
\end{tabular}
\end{center}

%======================================================================
\section{Data Pipeline}\label{sec:data}
%======================================================================

\subsection{Download and cleaning (Cell 5)}

Raw daily OHLCV comes from Yahoo Finance (\code{yfinance}, \code{auto\_adjust=False}).
Cleaning is deliberately conservative:

\begin{enumerate}
\item Forward-fill Close.  Backfill Open/High/Low from Close if missing.  Volume defaults to~0.
\item Drop any row that still has no Close.
\item Discard any ticker with $\le 100$ surviving daily bars.
\end{enumerate}

\subsection{Weekly resampling (Cell 6)}

Aggregation is anchored on Friday close:
$O=\text{first}$, $H=\max$, $L=\min$, $C=\text{last}$, $V=\text{sum}$.
Weekly returns and Garman--Klass volatility are computed per the
definitions in \textit{Math Foundations}, Section~1.

\subsection{SPY benchmark}

SPY is downloaded separately (\code{auto\_adjust=True}), resampled to the same
weekly grid, and aligned to the test-period index.  It is never part of the
tradable universe.

%======================================================================
\section{Feature Engineering (Cell 7)}\label{sec:features}
%======================================================================

Two feature sets are built, one for each prediction task.

\subsubsection*{Pair features (10-dim)}

One row per unordered pair $(i,j)$ per week.
Three lags of rolling correlation, both assets' returns and GK vols,
a vol ratio, and a cross-vol term.
Target: next-week correlation $\hat\rho_{ij,t+1}$.

\subsubsection*{Return features (14-dim)}

One row per ticker per week.
Five lagged returns, three rolling means (4/12/26\,w),
two rolling vols (8/12\,w), GK vol, vol-of-vol, and return acceleration.
Target: next-week return $r_{i,t+1}$.

Both feature sets are described equation-by-equation in
\textit{Math Foundations}, Section~3.

\placeholder{Insert feature dimension summary from execution output.}

%======================================================================
\section{Ensemble Architecture (Cell 8)}\label{sec:ensemble}
%======================================================================

\subsection{Base learners}

Nine models, loosely ordered from simple to complex:

\begin{center}
\renewcommand{\arraystretch}{1.15}
\begin{tabular}{@{}clcl@{}}
\toprule
\# & \textbf{Model} & \textbf{Tuned?} & \textbf{Key defaults}\\
\midrule
1 & OLS         & No  & ---\\
2 & Ridge       & Yes & $\alpha=1.0$; search $[10^{-4},100]$\\
3 & Lasso       & Yes & $\alpha=0.001$, \code{max\_iter}=5000\\
4 & ElasticNet  & Yes & $\alpha=0.001$, $\ell_1$-ratio$=0.5$\\
5 & Random Forest & Yes & 250 trees, depth 6, min-leaf 5\\
6 & Extra Trees & Yes & 250 trees, depth 8, min-leaf 3\\
7 & GBM         & Yes & 300 rounds, depth 3, $\eta=0.05$\\
8 & AdaBoost    & No  & 100 estimators, $\eta=0.05$\\
9 & MLP         & No  & 3$\times$64 hidden, ReLU, dropout 0.1\\
\bottomrule
\end{tabular}
\end{center}

Tuning uses Optuna (TPE sampler, seed~42) with 10 trials or 20\,s
wall-clock, whichever binds first.
CV is \code{TimeSeriesSplit(5)}, scored by MSE.

\subsection{MLP details}

\begin{lstlisting}[language=Python, caption={MLP architecture (identical for pair and return tasks).}]
layers = []
in_dim = d   # 10 (pair) or 14 (return)
for _ in range(3):
    layers += [Linear(in_dim, 64), ReLU(), Dropout(0.1)]
    in_dim = 64
layers.append(Linear(64, 1))
# Adam, lr=1e-3, weight_decay=1e-4, batch=64, patience=12
\end{lstlisting}

\subsection{Stacking}

After the base learners are trained on the first 80\% of in-sample data,
they predict on the held-out 20\% to produce a $2M$-dimensional meta-feature
vector: each base model's prediction, plus its absolute error.
A Ridge meta-model ($\alpha=1.0$) learns the final combination.

At test time, the absolute-error half is zeroed out
(the true target is unavailable), so the meta-model falls back
to a learned linear blend of the raw predictions.

\subsection{Training procedure}

For each pair or ticker:

\begin{enumerate}
\item Filter to training period.  Require $\ge 50$ rows.
\item 80/20 chronological split into base and meta sets.
\item \code{StandardScaler} fit on the base set, applied to both.
\item Train all 9 base learners (with Optuna where marked).
\item Predict on the meta set $\to$ meta-features.
\item Train Ridge meta-model.
\item Store scaler + all models.
\end{enumerate}

Per-model diagnostics (MSE, MAE, $R^2$, directional accuracy) are
written to \code{MODEL\_DIAG} for later plotting.

%======================================================================
\section{Prediction and Error Tracking (Cell 10)}\label{sec:predict}
%======================================================================

At each test date $t$, the trained ensembles produce:

\begin{itemize}
\item $\hat\mu_{t+1}\in\mathbb{R}^n$ (return predictions, unclipped).
\item $\hat C_{t+1}\in\mathbb{R}^{n\times n}$ (correlation predictions, clipped to $[-1,1]$).
\end{itemize}

Error histories grow as lists per ticker/pair.  At each step we append
the signed error $\varepsilon = \text{predicted} - \text{actual}$ and
recompute the five error statistics (EWMA bias, std, skewness, kurtosis,
directional hit) that feed the confidence module.
The EWMA uses $\lambda=0.94$; higher moments come from \code{scipy.stats}
on the full history.

%======================================================================
\section{Allocation Engine (Cell 11)}\label{sec:alloc}
%======================================================================

Five stages run in sequence for each test week.
The math is in \textit{Math Foundations}, Sections 4--9; here we
document the implementation.

\begin{tcolorbox}[colback=gray!3, colframe=gray!50!black]
\textbf{Stage 0: TC matrix} (\code{build\_tc\_matrix}).
Reads GK vol for the current week, computes $C_{ij}(t)$ per asset class.
$\to$ $C\in\mathbb{R}^{n\times n}$, symmetric, zero diagonal.

\textbf{Stage 1: Confidence} (\code{compute\_conf}).
Reads $E_\mu$, returns $\kappa\in\mathbb{R}^n$.

\textbf{Stage 2: Bayesian posterior} (\code{bayes\_post}).
NIW update $\to$ $(\mu^B,\Sigma^B)$.

\textbf{Stage 3: Topological analysis} (\code{topo\_analysis}).
Persistent homology on the correlation distance graph,
spectral denoising $\to$ $C^{\mathrm{spec}}$.

\textbf{Stage 4: Robust covariance} (\code{robust\_cov}).
EWMA vol from $V$, regime adjustment, blend with $\Sigma^B$
$\to$ $\Sigma_{\mathrm{final}}$.

\textbf{Stage 5: Sharpe allocation} (\code{sharpe\_alloc}).
50-point $\gamma$-sweep, CVXPY + SCS (\code{max\_iters}=5000),
pick $\gamma^*$ by Sharpe $\to$ $\bm{w}^*$.
\end{tcolorbox}

\textbf{Nearest-PD projection.}
Called after every covariance construction:
eigendecompose, clamp eigenvalues $\ge 10^{-8}$, reconstruct and symmetrise.
If the minimum eigenvalue is still below $10^{-8}$, add a diagonal ridge.

\textbf{Fallback.}
If any stage throws, the allocator returns equal-weight
($\bm{w}=n^{-1}\bm{1}$) and logs a warning.

%======================================================================
\section{Order Book Tracker (Cell 4)}\label{sec:ob}
%======================================================================

This is a simulation layer, not a live feed.  It reconstructs synthetic
L2 snapshots from weekly OHLCV so that the backtest can estimate
realistic spread and slippage costs.

\textbf{Snapshot construction.}
Mid = Friday close.
Spread $\approx \max(\hat\sigma^{\mathrm{GK}}\times 0.1,\;10^{-4})$.
Five depth levels with geometrically decaying size
($1000\times 0.8^\ell$, $\ell=0,\dots,4$).

\textbf{Rebalance logging.}
Each off-diagonal OT flow $T_{ij}^*>10^{-8}$ is logged as a fill,
recording: date, sell/buy assets, flow \%, notional, spread, slippage.
Weekly summaries (fill count, total notional, total slippage, turnover~\%)
feed the diagnostic plots.

%======================================================================
\section{Backtesting Engine (Cell 12)}\label{sec:backtest}
%======================================================================

One iteration per test week:

\begin{enumerate}
\item Retrieve actual weekly returns.
\item Build order book snapshot.
\item Run the five-stage allocation pipeline.
\item Compute OT transport plan and transaction cost $\mathcal{C}^*$.
\item Log the rebalance.
\item Raw return: $r^{\mathrm{raw}}=(\bm{w}^*)^\top\bm{r}-\mathcal{C}^*$.
\item Vol-target scaling (lookback 26\,w, cap $2.5\times$).
\item Record diagnostics: weights, $\kappa$, $\mu^B$, $\sigma$, leverage, turnover.
\end{enumerate}

%======================================================================
\section{Vol-Fitting and Metrics (Cell 13)}\label{sec:metrics}
%======================================================================

Three strategies --- PRISM, equal-weight, and SPY --- are all vol-fitted to
15\% annualised using the same mechanism (26-week lookback, cap $3.0\times$).
This gives a fair apples-to-apples comparison at identical risk budgets.

\textbf{Metrics computed} (per strategy):
total return, annualised return, annualised vol, Sharpe, Sortino,
max drawdown, Calmar, win rate, skewness, kurtosis.

%======================================================================
\section{Diagnostic Plots (Cell 14)}\label{sec:plots}
%======================================================================

32 plots, grouped thematically.

\begin{center}
\renewcommand{\arraystretch}{1.05}
\begin{tabular}{@{}rll@{}}
\toprule
& \textbf{Plot} & \textbf{Shows}\\
\midrule
\multicolumn{3}{@{}l}{\textit{Performance}}\\
P01 & Cumulative returns       & All strategies, log scale\\
P02 & Vol-matched comparison   & PRISM / EW / SPY at 15\%\\
P03 & Drawdown panels          & Three strategies stacked\\
P04 & Rolling Sharpe (26\,w)   & \\
P05 & Rolling return (52\,w)   & \\
P06 & Rolling vol + target     & 26\,w realised vs.\ 15\% line\\[4pt]
\multicolumn{3}{@{}l}{\textit{Portfolio composition}}\\
P07 & Weight stacks            & Long / short, stacked area\\
P08 & Leverage                 & Gross and net\\
P09 & Turnover + TC            & Weekly bars\\
P22 & Weight heatmap           & Ticker $\times$ time\\[4pt]
\multicolumn{3}{@{}l}{\textit{Return analysis}}\\
P10 & Return distributions     & Histogram overlay\\
P11 & QQ plots                 & vs.\ normal\\
P12 & Monthly heatmap          & Calendar grid\\
P23 & Annual returns           & Grouped bar chart\\
P24 & Underwater plot          & \\[4pt]
\multicolumn{3}{@{}l}{\textit{Model diagnostics}}\\
P13 & $\kappa$ evolution       & 10 representative tickers\\
P14 & Avg $\kappa$ + predicted vol & Dual-axis\\
P15 & $\mu^B$ heatmap          & \\
P16 & $\sigma$ heatmap         & \\
P17 & Final correlation + $\kappa$ & Last test date snapshot\\
P18 & Return $R^2$ per ticker  & \\
P19 & Return DA per ticker     & \\
P20 & Meta-model weights       & Ridge coefficients\\
P21 & Pair model diagnostics   & $R^2$, DA, MAE\\[4pt]
\multicolumn{3}{@{}l}{\textit{Market and execution}}\\
P25 & Average GK vol           & Cross-sectional\\
P26 & All assets (log)         & Normalised to 1\\
P27 & Individual asset panels  & Small multiples\\
P28 & Strategy vs.\ benchmarks & At 15\% vol\\
P29 & Order book summary       & Fill counts, notional\\
P30 & Spread history           & \\
P31 & Last trades              & Net flows, final rebalance\\
P32 & TC matrix heatmap        & Average transaction costs (bps)\\
\bottomrule
\end{tabular}
\end{center}

\end{document}